% ===============================================================
% Informe en formato APA 7.ª edición (trabajo estudiantil)
% Compilar con:  latexmk -pdfxe main.tex   (usa XeLaTeX + biber)
% ===============================================================
\documentclass[11pt,letterpaper]{article}

% --- Idioma y fuente -------------------------------------------
\usepackage{fontspec}
\usepackage[spanish,es-tabla]{babel}
\addto\captionsspanish{\renewcommand{\contentsname}{Índice}}

% Arial (APA 7 admite Arial 11 pt). Si no está instalada, usa
% Liberation Sans, que es métricamente idéntica.
\IfFontExistsTF{Arial}
  {\setmainfont{Arial}\setsansfont{Arial}}
  {\setmainfont{Liberation Sans}\setsansfont{Liberation Sans}}

% --- Datos de la portada (edita aquí) --------------------------
\newcommand{\titulo}{Aplicación de algoritmos de grafos en OSINT de archivos}
\newcommand{\estudiante}{Jhonatan Vargas Ramos}
\newcommand{\universidad}{Nombre de la universidad}
\newcommand{\carrera}{Ingeniería de Sistemas}
\newcommand{\asignatura}{Nombre de la asignatura}
\newcommand{\semestre}{Semestre}
\newcommand{\docente}{Nombre del docente}
\newcommand{\fecha}{\today}

% --- Página, interlineado y párrafos ---------------------------
\usepackage[margin=1in]{geometry}      % márgenes de 2,54 cm
\usepackage{setspace}
\doublespacing                         % doble espacio en todo el texto
\usepackage{ragged2e}
\RaggedRight                           % alineado a la izquierda
\setlength{\parindent}{0.5in}          % sangría de primera línea
\setlength{\parskip}{0pt}
\usepackage{indentfirst}               % también sangra tras un título
\setlength{\emergencystretch}{3em}     % evita Overfull \hbox con \RaggedRight
\usepackage{etoolbox}

% --- Número de página arriba a la derecha ----------------------
\usepackage{fancyhdr}
\pagestyle{fancy}
\fancyhf{}
\fancyhead[R]{\thepage}
\renewcommand{\headrulewidth}{0pt}
\setlength{\headheight}{14pt}
\fancypagestyle{plain}{\fancyhf{}\fancyhead[R]{\thepage}%
  \renewcommand{\headrulewidth}{0pt}}

% --- Niveles de título APA 7 -----------------------------------
\setcounter{secnumdepth}{0}            % títulos sin numeración
\setcounter{tocdepth}{3}
\usepackage{titlesec}
% Nivel 1: centrado, negrita
\titleformat{\section}{\normalfont\normalsize\bfseries\centering}{}{0pt}{}
% Nivel 2: izquierda, negrita
\titleformat{\subsection}{\normalfont\normalsize\bfseries}{}{0pt}{}
% Nivel 3: izquierda, negrita cursiva
\titleformat{\subsubsection}{\normalfont\normalsize\bfseries\itshape}{}{0pt}{}
% Nivel 4: sangrado, negrita, termina en punto, texto sigue en la línea
\titleformat{\paragraph}[runin]{\normalfont\normalsize\bfseries}{}{0pt}{}[.]
% Nivel 5: sangrado, negrita cursiva, termina en punto
\titleformat{\subparagraph}[runin]{\normalfont\normalsize\bfseries\itshape}{}{0pt}{}[.]
\titlespacing*{\section}{0pt}{0pt}{0pt}
\titlespacing*{\subsection}{0pt}{0pt}{0pt}
\titlespacing*{\subsubsection}{0pt}{0pt}{0pt}
\titlespacing*{\paragraph}{0.5in}{0pt}{1em}
\titlespacing*{\subparagraph}{0.5in}{0pt}{1em}

% --- Tablas y figuras APA --------------------------------------
\usepackage{booktabs}
\usepackage{graphicx}
\usepackage[labelfont=bf,textfont=it,labelsep=newline,
  singlelinecheck=false,justification=raggedright,
  position=top,skip=0pt]{caption}
\AtBeginEnvironment{table}{\singlespacing}
\AtBeginEnvironment{figure}{\singlespacing}
\newcommand{\nota}[1]{\par\vspace{0.5\baselineskip}\noindent\textit{Nota.} #1}

% --- Anexos ----------------------------------------------------
\newcounter{anexo}
\renewcommand{\theanexo}{\Alph{anexo}}
\newcommand{\anexo}[1]{%
  \ifnum\value{anexo}>0 \clearpage\fi
  \refstepcounter{anexo}%
  \begin{center}\textbf{Anexo \theanexo}\\ \textbf{#1}\end{center}%
  \addcontentsline{toc}{subsection}{Anexo \theanexo. #1}}

% --- Citas y referencias (APA 7 con biblatex-apa + biber) ------
\usepackage{csquotes}
\usepackage[backend=biber,style=apa]{biblatex}
\DeclareLanguageMapping{spanish}{spanish-apa}
\addbibresource{referencias.bib}
\defbibheading{bibliography}{\clearpage\section{#1}}

\usepackage[hidelinks,pdftitle={\titulo},pdfauthor={\estudiante}]{hyperref}
\usepackage{xurl}                       % permite cortar URLs y DOI largos

% ===============================================================
\begin{document}

% 1. Portada
% ---------------------------------------------------------------
% Portada estilo APA 7 (trabajo estudiantil)
% Usa las variables definidas en main.tex
% ---------------------------------------------------------------
\thispagestyle{fancy}
\vspace*{3\baselineskip}
\begin{center}
  \begin{center}\textbf{\titulo}\end{center}
  \vspace{1\baselineskip}

  \estudiante \\
  \carrera, \universidad \\
  \asignatura\ (\semestre) \\
  \docente \\
  \fecha
\end{center}
\clearpage


% 2. Índice
\begin{singlespace}
  \tableofcontents
\end{singlespace}
\clearpage

% En la primera página del texto, APA repite el título en negrita
\begin{center}\textbf{\titulo}\end{center}

\section{Introducción}
Escribe aquí el contenido de esta sección. Para citar usa
\parencite{cormen2022} o \textcite{dijkstra1959}.

\section{Descripción de la problemática}
Escribe aquí el contenido de esta sección.

\section{Objetivos}
\subsection{Objetivo general}
Escribe aquí el contenido de esta sección.
\subsection{Objetivos específicos}
Escribe aquí el contenido de esta sección.

\section{Metodología de recopilación de datos}
Escribe aquí el contenido de esta sección.

\section{Evidencias del contexto real}
Escribe aquí el contenido de esta sección.

\section{Fundamento teórico}
Escribe aquí el contenido de esta sección.

\section{Construcción y representación del grafo}
Escribe aquí el contenido de esta sección.

\section{Aplicación de DFS}
\subsection{Procedimiento}
Escribe aquí el contenido de esta sección.
\subsubsection{Resultados parciales}
Ejemplo de nivel 3.
\paragraph{Pila de exploración}
Ejemplo de nivel 4: el texto continúa en la misma línea del título.

\section{Aplicación de BFS}
Escribe aquí el contenido de esta sección.

\section{Aplicación de Dijkstra}
Escribe aquí el contenido de esta sección.

\section{Aplicación de Kruskal}
Escribe aquí el contenido de esta sección.

\section{Aplicación de Prim}
Escribe aquí el contenido de esta sección.

\section{Comparación de resultados}
La Tabla~\ref{tab:complejidad} resume la complejidad de cada algoritmo.

\begin{table}[ht]
  \caption{Complejidad temporal de los algoritmos aplicados}
  \label{tab:complejidad}
  \begin{tabular}{@{}lll@{}}
    \toprule
    Algoritmo & Estructura auxiliar & Complejidad \\
    \midrule
    DFS      & Pila             & $O(V+E)$ \\
    BFS      & Cola             & $O(V+E)$ \\
    Dijkstra & Cola de prioridad & $O((V+E)\log V)$ \\
    Kruskal  & Unión-búsqueda   & $O(E\log E)$ \\
    Prim     & Cola de prioridad & $O(E\log V)$ \\
    \bottomrule
  \end{tabular}
  \nota{$V$ = número de vértices; $E$ = número de aristas.}
\end{table}

% Ejemplo de figura (el título va ARRIBA de la imagen en APA):
% \begin{figure}[ht]
%   \caption{Grafo del caso de estudio}
%   \label{fig:grafo}
%   \includegraphics[width=0.8\textwidth]{grafo.png}
%   \nota{Elaboración propia.}
% \end{figure}

\section{Descripción del sistema}
Escribe aquí el contenido de esta sección.

\section{Casos de prueba}
Escribe aquí el contenido de esta sección.

\section{Interpretación de resultados}
Escribe aquí el contenido de esta sección.

\section{Recomendaciones}
Escribe aquí el contenido de esta sección.

\section{Conclusiones}
Escribe aquí el contenido de esta sección.

% 21. Referencias bibliográficas
\printbibliography[title={Referencias}]

% 22. Anexos
\clearpage
\section{Anexos}
\anexo{Código fuente del sistema}
Escribe aquí el contenido del anexo.

\anexo{Capturas de pantalla}
Escribe aquí el contenido del anexo.

\end{document}
